% ----- auto-generated from week_01/Week*.html via pandoc -----
\setchapterimage[6cm]{}
\setchapterpreamble[u]{\margintoc}
\chapter{Manipulator Dynamics and Computed-Torque Control (Week 1)}
\labch{week_01}



From Equations of Motion to Computed Torque Control

Control for Robots\\
AY 2025--26 \textbar{} 3-Hour Lecture + Activity Session

\section{Overview}

\subsection{Learning Outcomes}

\textbf{By the end of this session, you will be able to:}

\begin{itemize}
\tightlist
\item
  \textbf{Derive} the Euler--Lagrange equations of motion for a 2-DOF robot manipulator
\item
  \textbf{Identify} the mass matrix \(M(q)\), Coriolis matrix \(C(q,\dot{q})\), and gravity vector \(G(q)\) and explain their physical roles
\item
  \textbf{Design} a computed torque (inverse dynamics) controller and prove its stability using Lyapunov analysis
\item
  \textbf{Implement} computed torque control in MATLAB/Simulink and interpret tracking error plots
\item
  \textbf{Analyse} robustness of the controller against model uncertainties, payload variations, and disturbances
\item
  \textbf{Compare} PD, PD + gravity compensation, and full computed torque control performance
\end{itemize}

\subsection{Session Roadmap (3 Hours)}

0:00 -- 0:40

\subsubsection{Manipulator Dynamics}

\begin{itemize}
\tightlist
\item
  Review of feedback vs feedforward control
\item
  Euler--Lagrange formulation
\item
  Properties of the dynamics: \(M(q)\), \(C(q,\dot{q})\), \(G(q)\)
\item
  Skew-symmetry property \& its role in control
\end{itemize}

0:40 -- 1:20

\subsubsection{Control Design}

\begin{itemize}
\tightlist
\item
  PD control and its limitations
\item
  PD + gravity compensation
\item
  Computed torque (inverse dynamics) control
\item
  Lyapunov stability proof
\end{itemize}

1:20 -- 1:30

\subsubsection{Break}

\begin{itemize}
\tightlist
\item
  10-minute break
\end{itemize}

1:30 -- 2:15

\subsubsection{MATLAB Implementation}

\begin{itemize}
\tightlist
\item
  Guided MATLAB coding of computed torque controller
\item
  Running ODE45 simulations
\item
  Interpreting tracking and error plots
\end{itemize}

2:15 -- 3:00

\subsubsection{Robustness Analysis \& Activities}

\begin{itemize}
\tightlist
\item
  Effect of model uncertainty on computed torque
\item
  Payload variation experiments
\item
  Disturbance rejection analysis
\item
  Group activity \& wrap-up quiz
\end{itemize}

\section{Manipulator Dynamics}

\subsection{Why Dynamics Matter}

Kinematics tells us \emph{where} the robot moves; dynamics tells us \emph{what forces and torques} are required to produce that motion. Without a dynamic model we cannot design controllers that truly cancel the nonlinear effects of inertia, coupling, Coriolis forces, and gravity.

\subsubsection{Feedback Control}

Uses sensor measurements (encoders, IMUs) to adjust torques in real time. Robust against model errors but reacts \emph{after} errors occur.

\subsubsection{Feedforward Control}

Applies precomputed torques from the dynamic model \emph{before} the error appears. Requires an accurate model but gives faster response.

\subsubsection{Combined (Computed Torque)}

Feedforward cancels known dynamics; feedback corrects remaining error. This is the gold standard for manipulator control.

\subsection{The Euler--Lagrange Equation of Motion}

For an \(n\)-DOF rigid-body manipulator, the Lagrangian \(\mathcal{L} = T - V\) (kinetic minus potential energy) yields the standard form:

\[ M(q)\,\ddot{q} \;+\; C(q,\dot{q})\,\dot{q} \;+\; G(q) \;=\; \tau \]

\begin{longtable}[]{@{}llll@{}}
\toprule\noalign{}
Symbol & Name & Dimension & Physical Meaning \\
\midrule\noalign{}
\endhead
\bottomrule\noalign{}
\endlastfoot
\(q\) & Joint positions & \(n \times 1\) & Generalised coordinates (angles for revolute joints) \\
\(M(q)\) & Mass (inertia) matrix & \(n \times n\) & Relates joint accelerations to required torques; symmetric positive-definite \\
\(C(q,\dot{q})\) & Coriolis \& centrifugal matrix & \(n \times n\) & Velocity-dependent coupling forces between joints \\
\(G(q)\) & Gravity vector & \(n \times 1\) & Torques required to hold the arm against gravity \\
\(\tau\) & Joint torques & \(n \times 1\) & Control inputs (what we design) \\
\end{longtable}

\subsection{A Concrete Example: 2-DOF Planar Manipulator}

Consider a two-link planar arm with link masses \(m_1, m_2\), link lengths \(l_1, l_2\), and distances to centres of mass \(l_{c1}, l_{c2}\). For simplicity we take \(l_{ci} = l_i\) (point masses at link tips).

\paragraph{\texorpdfstring{📐 Mass Matrix \(M(q)\)}{📐 Mass Matrix M(q)}}}

\[
                    M(q) = \begin{bmatrix}
                    m_1 l_1^2 + m_2(l_1^2 + l_2^2 + 2l_1 l_2 \cos q_2) & m_2(l_2^2 + l_1 l_2 \cos q_2) \\
                    m_2(l_2^2 + l_1 l_2 \cos q_2) & m_2 l_2^2
                    \end{bmatrix}
                    \]

Notice: \(M(q)\) depends on \(q_2\) only (the relative angle), and is always symmetric positive-definite so it is invertible.

\paragraph{\texorpdfstring{🌀 Coriolis \& Centrifugal Matrix \(C(q,\dot{q})\)}{🌀 Coriolis \& Centrifugal Matrix C(q,\textbackslash dot\{q\})}}}

Using Christoffel symbols of the first kind:

\[
                    C(q,\dot{q}) = \begin{bmatrix}
                    -m_2 l_1 l_2 \sin(q_2)\,\dot{q}_2 & -m_2 l_1 l_2 \sin(q_2)(\dot{q}_1 + \dot{q}_2) \\
                    m_2 l_1 l_2 \sin(q_2)\,\dot{q}_1 & 0
                    \end{bmatrix}
                    \]

\paragraph{\texorpdfstring{⬇️ Gravity Vector \(G(q)\)}{⬇️ Gravity Vector G(q)}}}

\[
                    G(q) = \begin{bmatrix}
                    (m_1 + m_2)\,g\,l_1 \cos q_1 + m_2\,g\,l_2 \cos(q_1 + q_2) \\
                    m_2\,g\,l_2 \cos(q_1 + q_2)
                    \end{bmatrix}
                    \]

\subsection{Key Properties of the Dynamic Model}

\subsubsection{Property 1: Symmetry \& Positive-Definiteness}

\(M(q) = M(q)^T > 0\) for all \(q\). This guarantees the system always has finite inertia and the dynamics are well-posed.

\subsubsection{Property 2: Skew-Symmetry}

The matrix \(\dot{M}(q) - 2C(q,\dot{q})\) is skew-symmetric, i.e., \(x^T[\dot{M} - 2C]x = 0\) for all \(x\). This property is \textbf{essential} for Lyapunov-based stability proofs.

\subsubsection{Property 3: Linearity in Parameters}

The dynamics can be written as \(M(q)\ddot{q} + C(q,\dot{q})\dot{q} + G(q) = Y(q,\dot{q},\ddot{q})\,\theta\) where \(\theta\) contains the physical parameters. This enables adaptive control.

\subsubsection{Think--Pair--Share (5 min)}

\textbf{Question:} If the second link mass \(m_2\) doubles (e.g., the robot picks up a heavy payload), which terms in \(M(q)\), \(C(q,\dot{q})\), and \(G(q)\) are affected? Which are \emph{not} affected?

Discuss with your neighbour for 2 minutes, then we will share answers as a class.

\section{Control Design for the Manipulator}

\subsection{Approach 1: Simple PD Control}

The simplest approach is a joint-space PD controller:

\[ \tau = K_P\,e + K_D\,\dot{e} \]

where \(e = q_d - q\) is the tracking error and \(K_P, K_D \succ 0\) are diagonal gain matrices.

This treats the manipulator as if it were a collection of independent linear joints. It ignores nonlinear coupling, Coriolis forces, and gravity entirely.

\textbf{Limitation:} PD control alone cannot guarantee zero steady-state error in the presence of gravity. The arm will droop under its own weight because no term compensates for \(G(q)\).

\subsection{Approach 2: PD + Gravity Compensation}

\[ \tau = K_P\,e + K_D\,\dot{e} + G(q) \]

Adding the gravity term \(G(q)\) eliminates the steady-state droop at the desired position. However, during fast motions the Coriolis and inertia coupling terms still cause significant tracking errors.

\subsection{Approach 3: Computed Torque Control (Inverse Dynamics)}

The computed torque controller uses the full dynamic model to \textbf{cancel all nonlinear dynamics} and then imposes a desired linear closed-loop behaviour.

\textbf{Computed Torque Control Law}\\
\strut \\
\[ \tau = M(q)\,\bigl[\ddot{q}_d + K_D\,\dot{e} + K_P\,e\bigr] + C(q,\dot{q})\,\dot{q} + G(q) \]

\subsubsection{Derivation --- Step by Step}

\textbf{Step 1:} Start from the robot dynamics:

\[ M(q)\,\ddot{q} + C(q,\dot{q})\,\dot{q} + G(q) = \tau \]

\textbf{Step 2:} Substitute the control law into the dynamics:

\[ M(q)\,\ddot{q} + \cancel{C(q,\dot{q})\,\dot{q}} + \cancel{G(q)} = M(q)\bigl[\ddot{q}_d + K_D\dot{e} + K_P e\bigr] + \cancel{C(q,\dot{q})\,\dot{q}} + \cancel{G(q)} \]

\textbf{Step 3:} Since \(M(q)\) is invertible, left-multiply both sides by \(M^{-1}(q)\):

\[ \ddot{q} = \ddot{q}_d + K_D\,\dot{e} + K_P\,e \]

\textbf{Step 4:} Since \(\ddot{e} = \ddot{q}_d - \ddot{q}\), rearrange to get the \textbf{closed-loop error dynamics}:

\[ \ddot{e} + K_D\,\dot{e} + K_P\,e = 0 \]

This is a \textbf{linear, decoupled, second-order ODE} for each joint --- identical to a mass--spring--damper system. With proper gain selection the error converges to zero exponentially!

\subsection{Gain Selection: Placing Closed-Loop Poles}

For each joint, the characteristic equation is \(s^2 + k_d\,s + k_p = 0\). Comparing with the standard second-order form \(s^2 + 2\zeta\omega_n s + \omega_n^2 = 0\):

\[ k_p = \omega_n^2, \qquad k_d = 2\zeta\omega_n \]

\begin{longtable}[]{@{}llllll@{}}
\toprule\noalign{}
Design Choice & \(\zeta\) & \(\omega_n\) & \(K_P\) & \(K_D\) & Behaviour \\
\midrule\noalign{}
\endhead
\bottomrule\noalign{}
\endlastfoot
Underdamped & 0.5 & 10 & 100 & 10 & Fast but oscillatory \\
Critically damped & 1.0 & 10 & 100 & 20 & Fastest without overshoot ✓ \\
Overdamped & 2.0 & 10 & 100 & 40 & Sluggish but smooth \\
\end{longtable}

\subsection{Lyapunov Stability Proof}

Even though the computed torque controller linearises the system, we can also prove stability directly using Lyapunov theory --- which is more rigorous and generalises to uncertain models.

\subsubsection{Lyapunov Proof for Computed Torque Control}

\textbf{Step 1 --- Choose a Lyapunov candidate:}

\[ V(e, \dot{e}) = \frac{1}{2}\dot{e}^T M(q)\,\dot{e} + \frac{1}{2}e^T K_P\,e \]

Since \(M(q) \succ 0\) and \(K_P \succ 0\), we have \(V > 0\) for all \((e, \dot{e}) \ne (0,0)\).

\textbf{Step 2 --- Differentiate:}

\[ \dot{V} = \dot{e}^T M(q)\,\ddot{e} + \frac{1}{2}\dot{e}^T \dot{M}(q)\,\dot{e} + e^T K_P\,\dot{e} \]

\textbf{Step 3 --- Substitute the closed-loop dynamics} \(M(q)\ddot{e} = -C(q,\dot{q})\dot{e} - K_D\dot{e} - K_P e\):

\[ \dot{V} = \dot{e}^T\bigl[-C\dot{e} - K_D\dot{e} - K_P e\bigr] + \frac{1}{2}\dot{e}^T\dot{M}\,\dot{e} + e^T K_P\dot{e} \]

\textbf{Step 4 --- Apply the skew-symmetry property} \(\dot{e}^T[\dot{M} - 2C]\dot{e} = 0\):

\[ \dot{V} = -\dot{e}^T K_D\,\dot{e} \;\leq\; 0 \]

\textbf{Step 5 --- Invoke LaSalle\textquotesingle s invariance principle:} \(\dot{V} = 0 \implies \dot{e} = 0 \implies \ddot{e} = 0 \implies K_P e = 0 \implies e = 0\). Therefore the equilibrium \((e, \dot{e}) = (0, 0)\) is \textbf{asymptotically stable}. ∎

\textbf{Key Insight:} The skew-symmetry of \(\dot{M} - 2C\) is what makes the cross terms cancel in the Lyapunov derivative. This is the most important structural property in robot control!

\subsubsection{Quick Check (2 min)}

\begin{enumerate}
\tightlist
\item
  What happens if \(K_P = 0\) in the Lyapunov function? Can we still guarantee that \(e \to 0\)?
\item
  In the computed torque law, which term handles gravity? Which handles Coriolis coupling?
\item
  If \(M(q)\) were identity (unit mass everywhere), how would the computed torque law simplify?
\end{enumerate}

\section{MATLAB Implementation}

We will now implement the computed torque controller for a 2-DOF manipulator step by step. The robot must track a sinusoidal trajectory:
\(q_d(t) = \begin{bmatrix} 0.5\sin(t) \\ 0.5\cos(t) \end{bmatrix}\).

\subsection{System Parameters}

\% ===== System Parameters =====
clear all; close all; clc;
params.m1 = 1; \% Mass of link 1 (kg)
params.m2 = 1; \% Mass of link 2 (kg)
params.l1 = 1; \% Length of link 1 (m)
params.l2 = 1; \% Length of link 2 (m)
params.g = 9.81; \% Gravitational acceleration (m/s\^{}2)
\% Control gains (critically damped, wn = 7.07)
params.Kp = diag({[}50 50{]}); \% Kp = wn\^{}2
params.Kd = diag({[}20 20{]}); \% Kd = 2*zeta*wn (zeta ≈ 1.4)
\% Simulation setup
t\_span = 0:0.01:20; \% 20 seconds, 10 ms step
q0 = {[}0; 0{]}; \% Initial joint angles
dq0 = {[}0; 0{]}; \% Initial joint velocities
x0 = {[}q0; dq0{]}; \% State vector {[}q1; q2; dq1; dq2{]}

\subsection{Desired Trajectory Functions}

\% ===== Desired Trajectory =====
qd = @(t) {[}0.5*sin(t); 0.5*cos(t){]};
dqd = @(t) {[}0.5*cos(t); -0.5*sin(t){]};
ddqd = @(t) {[}-0.5*sin(t); -0.5*cos(t){]};

\subsection{Dynamics with Computed Torque Controller}

\subsubsection{Pseudocode}

📋 PSEUDOCODE --- Computed Torque Controller

FUNCTION manipulator\_dynamics(t, state, params):
EXTRACT q = state{[}1:2{]}, dq = state{[}3:4{]}
COMPUTE desired trajectory: qd(t), dqd(t), ddqd(t)
COMPUTE tracking error: e = qd - q
COMPUTE velocity error: de = dqd - dq
BUILD M(q): 2×2 mass matrix using m1, m2, l1, l2, cos(q2)
BUILD C(q,dq): 2×2 Coriolis matrix using m2, l1, l2, sin(q2), dq
BUILD G(q): 2×1 gravity vector using g, m1, m2, l1, l2
COMPUTE feedforward: a\_cmd = ddqd + Kd*de + Kp*e
COMPUTE torque: tau = M * a\_cmd + C * dq + G
COMPUTE acceleration: ddq = M \textbackslash{} (tau - C*dq - G)
RETURN {[}dq; ddq{]}

\subsubsection{MATLAB Code}

function dx = manipulator\_dynamics(t, x, params, qd, dqd, ddqd)
q = x(1:2); \% Joint angles
dq = x(3:4); \% Joint velocities
\% Desired trajectory at current time
q\_des = qd(t);
dq\_des = dqd(t);
ddq\_des = ddqd(t);
\% Tracking errors
e = q\_des - q;
de = dq\_des - dq;
\% -\/-\/- Build the Dynamic Model -\/-\/-
m1 = params.m1; m2 = params.m2;
l1 = params.l1; l2 = params.l2;
g = params.g;
\% Mass matrix M(q)
M11 = m1*l1\^{}2 + m2*(l1\^{}2 + l2\^{}2 + 2*l1*l2*cos(q(2)));
M12 = m2*(l2\^{}2 + l1*l2*cos(q(2)));
M22 = m2*l2\^{}2;
M = {[}M11, M12; M12, M22{]};
\% Coriolis matrix C(q, dq)
h = m2*l1*l2*sin(q(2));
C = {[}-h*dq(2), -h*(dq(1)+dq(2));
h*dq(1), 0 {]};
\% Gravity vector G(q)
G = {[}(m1+m2)*g*l1*cos(q(1)) + m2*g*l2*cos(q(1)+q(2));
m2*g*l2*cos(q(1)+q(2)) {]};
\% -\/-\/- Computed Torque Control Law -\/-\/-
a\_cmd = ddq\_des + params.Kd*de + params.Kp*e;
tau = M * a\_cmd + C * dq + G;
\% -\/-\/- Forward Dynamics -\/-\/-
ddq = M \textbackslash{} (tau - C*dq - G);
dx = {[}dq; ddq{]};
end

\textbf{Why \texttt{M\textbackslash{}b} instead of \texttt{inv(M)*b}?} The backslash operator is numerically more stable and faster. Never invert the mass matrix explicitly in production code.

\subsection{Run Simulation \& Plot}

\% ===== Run ODE45 =====
{[}t, x{]} = ode45(@(t,x) manipulator\_dynamics(t, x, params, qd, dqd, ddqd), ...
t\_span, x0);
\% ===== Compute desired trajectory for plotting =====
qd\_vals = zeros(length(t), 2);
for i = 1:length(t)
qd\_vals(i,:) = qd(t(i))\textquotesingle;
end
\% ===== Plot: Joint Angles =====
figure(\textquotesingle Position\textquotesingle, {[}100 100 1200 500{]});
subplot(1,2,1);
plot(t, x(:,1), \textquotesingle b\textquotesingle, \textquotesingle LineWidth\textquotesingle, 1.5); hold on;
plot(t, x(:,2), \textquotesingle r\textquotesingle, \textquotesingle LineWidth\textquotesingle, 1.5);
plot(t, qd\_vals(:,1), \textquotesingle b-\/-\textquotesingle, \textquotesingle LineWidth\textquotesingle, 1.2);
plot(t, qd\_vals(:,2), \textquotesingle r-\/-\textquotesingle, \textquotesingle LineWidth\textquotesingle, 1.2);
xlabel(\textquotesingle Time (s)\textquotesingle); ylabel(\textquotesingle Angle (rad)\textquotesingle);
title(\textquotesingle Joint Angles: Actual vs Desired\textquotesingle);
legend(\textquotesingle q\_1\textquotesingle, \textquotesingle q\_2\textquotesingle, \textquotesingle q\_\{1,d\}\textquotesingle, \textquotesingle q\_\{2,d\}\textquotesingle);
grid on;
\% ===== Plot: Tracking Error =====
subplot(1,2,2);
e\_vals = qd\_vals - x(:,1:2);
plot(t, e\_vals(:,1), \textquotesingle b\textquotesingle, \textquotesingle LineWidth\textquotesingle, 1.5); hold on;
plot(t, e\_vals(:,2), \textquotesingle r\textquotesingle, \textquotesingle LineWidth\textquotesingle, 1.5);
xlabel(\textquotesingle Time (s)\textquotesingle); ylabel(\textquotesingle Error (rad)\textquotesingle);
title(\textquotesingle Tracking Error\textquotesingle);
legend(\textquotesingle e\_1\textquotesingle, \textquotesingle e\_2\textquotesingle);
grid on;
sgtitle(\textquotesingle Computed Torque Control --- 2-DOF Manipulator\textquotesingle);

\textbf{Expected Result:} After a brief transient (≈ 1--2 s), both joint errors should converge to near zero and the actual trajectory should overlay the desired trajectory almost perfectly.

\subsubsection{Hands-On Activity 1 (10 min)}

\textbf{Task:} Run the code above in MATLAB. Then modify the gains as follows and observe the effect on the tracking error plot:

\begin{enumerate}
\tightlist
\item
  Set \(K_P = \text{diag}([10, 10])\) and \(K_D = \text{diag}([5, 5])\) --- \emph{what happens?}
\item
  Set \(K_P = \text{diag}([200, 200])\) and \(K_D = \text{diag}([5, 5])\) --- \emph{is there overshoot?}
\item
  Set \(K_P = \text{diag}([50, 50])\) and \(K_D = \text{diag}([50, 50])\) --- \emph{is the response sluggish?}
\end{enumerate}

Record the peak error for each case. We will compare results as a class.

\section{Comparing Control Strategies}

To appreciate why computed torque is superior, let us compare three controllers on the same trajectory.

\subsection{Controller A: PD Only}

\% Inside dynamics function, replace control law with:
tau = params.Kp * e + params.Kd * de;
\% No model terms --- ignores M, C, G completely

\subsection{Controller B: PD + Gravity Compensation}

\% Inside dynamics function, replace control law with:
tau = params.Kp * e + params.Kd * de + G;
\% Compensates gravity but ignores Coriolis and inertia coupling

\subsection{Controller C: Computed Torque (Full)}

\% Full computed torque (already implemented above):
a\_cmd = ddq\_des + params.Kd * de + params.Kp * e;
tau = M * a\_cmd + C * dq + G;

\begin{longtable}[]{@{}lllll@{}}
\toprule\noalign{}
Controller & Model Knowledge & Steady-State Error & Transient Error & Coupling Rejection \\
\midrule\noalign{}
\endhead
\bottomrule\noalign{}
\endlastfoot
PD Only & None & Non-zero (gravity droop) & Large & Poor \\
PD + Gravity & \(G(q)\) & Zero at setpoints & Moderate & Moderate \\
Computed Torque & \(M, C, G\) & Zero & Small (tunable) & Excellent ✓ \\
\end{longtable}

\subsubsection{Hands-On Activity 2 (15 min)}

\textbf{Task:} Implement all three controllers. Plot their tracking errors on the same figure using \texttt{subplot} or overlaid with different colours. Answer:

\begin{enumerate}
\tightlist
\item
  Which controller has the largest peak error during the first 2 seconds?
\item
  Which controller has a non-zero steady-state error? Why?
\item
  At what joint velocity does the PD controller start to diverge noticeably from computed torque?
\end{enumerate}

\section{Robustness Analysis}

Computed torque control delivers \textbf{perfect} tracking when the model is exact. In practice, we never know the model perfectly. Robustness analysis studies what happens when the model used in the controller differs from the true plant.

\subsection{Sources of Model Uncertainty}

\subsubsection{Payload Variation}

The mass \(m_2\) changes when the robot picks up or drops objects. The controller may still use the nominal value.

\subsubsection{Friction}

Real joints exhibit Coulomb and viscous friction not captured in the rigid-body model.

\subsubsection{Parameter Errors}

Link lengths and centres of mass are measured imprecisely. Actuator dynamics are neglected.

\subsubsection{External Disturbances}

Contact forces, wind loads, or interaction forces when collaborating with humans.

\subsection{Mathematical Framework}

Let the true dynamics be governed by \((M, C, G)\) and the \emph{estimated} (controller) model be \((\hat{M}, \hat{C}, \hat{G})\). The control law uses the estimated model:

\[ \tau = \hat{M}(q)\bigl[\ddot{q}_d + K_D\dot{e} + K_P e\bigr] + \hat{C}(q,\dot{q})\dot{q} + \hat{G}(q) \]

Substituting into the true dynamics:

\[ M\,\ddot{q} + C\,\dot{q} + G = \hat{M}\bigl[\ddot{q}_d + K_D\dot{e} + K_P e\bigr] + \hat{C}\,\dot{q} + \hat{G} \]

Defining model errors \(\Delta M = M - \hat{M}\), \(\Delta C = C - \hat{C}\), \(\Delta G = G - \hat{G}\), the closed-loop error dynamics become:

\[ M\,\ddot{e} + (C + K_D^*)\,\dot{e} + K_P^*\,e = \underbrace{\Delta M\,\ddot{q}_d + \Delta C\,\dot{q} + \Delta G}_{\text{perturbation } d(t)} \]

The error no longer converges to zero --- it is driven by the perturbation \(d(t)\).

\textbf{Key Result:} With model mismatch, computed torque control guarantees \emph{bounded} errors (the errors stay small if the gains are high and the model errors are small), but not \emph{asymptotic} tracking to zero.

\subsection{Experiment: Payload Mismatch}

We simulate the scenario where the \textbf{true} plant has \(m_2 = 2\,\text{kg}\) but the controller still assumes \(m_2 = 1\,\text{kg}\).

\% ===== Robustness Test: Payload Mismatch =====
\% True plant parameters (what reality is)
params\_true = params;
params\_true.m2 = 2; \% True mass of link 2 is DOUBLE
\% Controller still uses params (m2 = 1) for its model
\% We need a modified dynamics function:
function dx = manipulator\_robust\_test(t, x, params\_ctrl, params\_true, qd, dqd, ddqd)
q = x(1:2);
dq = x(3:4);
q\_des = qd(t);
dq\_des = dqd(t);
ddq\_des = ddqd(t);
e = q\_des - q;
de = dq\_des - dq;
\% -\/-\/- Controller uses ESTIMATED model (params\_ctrl) -\/-\/-
{[}M\_hat, C\_hat, G\_hat{]} = build\_model(q, dq, params\_ctrl);
a\_cmd = ddq\_des + params\_ctrl.Kd*de + params\_ctrl.Kp*e;
tau = M\_hat * a\_cmd + C\_hat * dq + G\_hat;
\% -\/-\/- Plant uses TRUE model (params\_true) -\/-\/-
{[}M\_true, C\_true, G\_true{]} = build\_model(q, dq, params\_true);
ddq = M\_true \textbackslash{} (tau - C\_true*dq - G\_true);
dx = {[}dq; ddq{]};
end
function {[}M, C, G{]} = build\_model(q, dq, p)
M11 = p.m1*p.l1\^{}2 + p.m2*(p.l1\^{}2 + p.l2\^{}2 + 2*p.l1*p.l2*cos(q(2)));
M12 = p.m2*(p.l2\^{}2 + p.l1*p.l2*cos(q(2)));
M22 = p.m2*p.l2\^{}2;
M = {[}M11, M12; M12, M22{]};
h = p.m2*p.l1*p.l2*sin(q(2));
C = {[}-h*dq(2), -h*(dq(1)+dq(2)); h*dq(1), 0{]};
G = {[}(p.m1+p.m2)*p.g*p.l1*cos(q(1)) + p.m2*p.g*p.l2*cos(q(1)+q(2));
p.m2*p.g*p.l2*cos(q(1)+q(2)){]};
end

\subsection{Experiment: Adding Joint Friction}

\% ===== Add Coulomb + Viscous Friction =====
\% After computing ddq from the model, add friction terms:
b\_viscous = 0.5; \% Viscous friction coefficient
f\_coulomb = 0.3; \% Coulomb friction magnitude
\% Friction torque (not modelled by controller)
tau\_friction = b\_viscous * dq + f\_coulomb * sign(dq);
\% Modified acceleration (true plant sees friction)
ddq = M\_true \textbackslash{} (tau - C\_true*dq - G\_true - tau\_friction);

\subsection{Experiment: External Disturbance Pulse}

\% ===== External Disturbance =====
\% Simulate a sudden impact at t = 10 s lasting 0.5 s
if t \textgreater= 10 \&\& t \textless= 10.5
tau\_dist = {[}5; 3{]}; \% Disturbance torque (N·m)
else
tau\_dist = {[}0; 0{]};
end
ddq = M\_true \textbackslash{} (tau - C\_true*dq - G\_true + tau\_dist);

\subsection{Summary of Robustness Results}

\begin{longtable}[]{@{}llll@{}}
\toprule\noalign{}
Scenario & Peak Error (rad) & Steady-State Error & Stability \\
\midrule\noalign{}
\endhead
\bottomrule\noalign{}
\endlastfoot
Perfect model & ≈ 0.02 (transient only) & ≈ 0 & Asymptotically stable \\
Payload ×2 & ≈ 0.08 -- 0.15 & Small oscillatory & Bounded (stable) \\
Joint friction & ≈ 0.03 -- 0.06 & Small bias & Bounded (stable) \\
Disturbance pulse & ≈ 0.10 (at impact) & Returns to ≈ 0 & ISS (input-to-state stable) \\
Payload ×5 + friction & ≈ 0.3 -- 0.5 & Significant & May become unstable with low gains \\
\end{longtable}

\textbf{Rule of Thumb:} Higher gains (\(K_P, K_D\)) improve robustness by rejecting perturbations faster, but they also amplify sensor noise and demand higher torques. Controller design is always a trade-off between performance, robustness, and actuator limits.

\section{Summary and Next Steps}

\subsection{Hands-On Activities}

\subsection{Activity 3: Robustness Investigation (20 min)}

\begin{enumerate}
\tightlist
\item
  Start with the nominal computed torque controller (\(m_2 = 1\,\text{kg}\) in both plant and controller).
\item
  Change the \textbf{true} \(m_2\) to 1.5, 2.0, 3.0, and 5.0 kg while keeping the controller model at \(m_2 = 1\).
\item
  For each case, record the maximum tracking error \(\|e\|_\infty\) from the error plot.
\item
  Plot \(\|e\|_\infty\) vs payload ratio \(m_{2,\text{true}} / m_{2,\text{ctrl}}\).
\item
  \textbf{Bonus:} At what payload ratio does the system become unstable? Increase \(K_P\) and \(K_D\) --- does this extend the stability boundary?
\end{enumerate}

\textbf{Expected Observation:} Tracking error grows approximately linearly with payload mismatch for small mismatches, but growth accelerates for large mismatches.

\subsection{Activity 4: Controller Comparison Dashboard (25 min)}

Create a single MATLAB figure with a 3×2 subplot layout:

\begin{longtable}[]{@{}lll@{}}
\toprule\noalign{}
& Column 1: Joint Angles & Column 2: Tracking Error \\
\midrule\noalign{}
\endhead
\bottomrule\noalign{}
\endlastfoot
\textbf{Row 1} & PD Only --- actual vs desired & PD Only --- error \\
\textbf{Row 2} & PD + Gravity --- actual vs desired & PD + Gravity --- error \\
\textbf{Row 3} & Computed Torque --- actual vs desired & Computed Torque --- error \\
\end{longtable}

Use the same gains and trajectory for all three. Include a title, legends, and axis labels on every subplot.

\subsection{Activity 5: Design Challenge (15 min)}

\subsubsection{Challenge: Robust Trajectory Tracking}

Design a controller that can track the reference trajectory with \(\|e\|_\infty < 0.05\,\text{rad}\) under the following \textbf{simultaneous} conditions:

\begin{itemize}
\tightlist
\item
  True \(m_2 = 2\,\text{kg}\), controller assumes \(m_2 = 1\,\text{kg}\)
\item
  Viscous friction \(b = 0.5\,\text{N·m·s/rad}\) (not modelled)
\item
  An external disturbance of \(\tau_d = [3; 2]\,\text{N·m}\) applied from \(t = 8\) to \(t = 8.5\,\text{s}\)
\end{itemize}

\textbf{Hint:} You may increase gains, add integral action, or add a robust term \(\tau_{\text{robust}} = K_r\,\text{sign}(\dot{e})\).

\subsection{End-of-Lecture Quiz}

\subsubsection{Test Your Understanding}

\begin{enumerate}
\item
  \textbf{True or False:} The mass matrix \(M(q)\) is always diagonal for a 2-DOF manipulator.

  \emph{Think about what the off-diagonal terms represent physically.}
\item
  \textbf{Fill in the blank:} The property that \(\dot{M} - 2C\) is \_\_\_\_\_\_\_\_\_\_\_ is essential for the Lyapunov stability proof of computed torque control.
\item
  \textbf{Short answer:} Name two practical reasons why computed torque control may not achieve perfect tracking on a real robot.
\item
  \textbf{Calculation:} For a desired natural frequency \(\omega_n = 15\,\text{rad/s}\) and critical damping (\(\zeta = 1\)), compute \(K_P\) and \(K_D\) per joint.
\item
  \textbf{Design question:} If you double all link masses simultaneously, but the controller model remains unchanged, how does the perturbation term \(d(t)\) scale? Would increasing \(K_P\) alone fix the problem? Why or why not?
\end{enumerate}

Click to reveal answers

\begin{enumerate}
\tightlist
\item
  \textbf{False.} The off-diagonal terms \(M_{12} = M_{21}\) capture inertia coupling between joints. They are zero only in special configurations.
\item
  \textbf{Skew-symmetric.}
\item
  Any two of: model parameter uncertainty, unmodelled friction, actuator saturation/dynamics, sensor noise, discretisation effects, joint flexibility.
\item
  \(K_P = \omega_n^2 = 225\), \(K_D = 2 \times 1 \times 15 = 30\).
\item
  The perturbation scales proportionally with the mass error \(\Delta M\), which doubles. Increasing \(K_P\) alone gives faster error correction but does not cancel the feedforward mismatch --- you also need higher \(K_D\) to damp the resulting oscillations, or better yet, update the model.
\end{enumerate}

\subsection{Week 1 Summary}

\subsubsection{Manipulator Dynamics}

\begin{itemize}
\tightlist
\item
  ✅ Euler--Lagrange formulation
\item
  ✅ Mass matrix \(M(q)\) --- SPD
\item
  ✅ Coriolis matrix \(C(q,\dot{q})\)
\item
  ✅ Gravity vector \(G(q)\)
\item
  ✅ Skew-symmetry property
\end{itemize}

\subsubsection{Control Design}

\begin{itemize}
\tightlist
\item
  ✅ PD control \& limitations
\item
  ✅ PD + gravity compensation
\item
  ✅ Computed torque derivation
\item
  ✅ Gain selection via pole placement
\item
  ✅ Lyapunov stability proof
\end{itemize}

\subsubsection{Implementation}

\begin{itemize}
\tightlist
\item
  ✅ MATLAB ODE45 simulation
\item
  ✅ Building \(M\), \(C\), \(G\) in code
\item
  ✅ Tracking \& error plots
\item
  ✅ Controller comparison
\item
  ✅ Gain tuning experiments
\end{itemize}

\subsubsection{Robustness}

\begin{itemize}
\tightlist
\item
  ✅ Payload mismatch analysis
\item
  ✅ Unmodelled friction effects
\item
  ✅ Disturbance rejection
\item
  ✅ Bounded error guarantees
\item
  ✅ Gain--robustness trade-off
\end{itemize}

\subsection{Key Equations to Remember}

\begin{longtable}[]{@{}ll@{}}
\toprule\noalign{}
Equation & Description \\
\midrule\noalign{}
\endhead
\bottomrule\noalign{}
\endlastfoot
\(M(q)\ddot{q} + C(q,\dot{q})\dot{q} + G(q) = \tau\) & Robot dynamics (Euler--Lagrange) \\
\(\tau = M(q)[\ddot{q}_d + K_D\dot{e} + K_P e] + C\dot{q} + G\) & Computed torque control law \\
\(\ddot{e} + K_D\dot{e} + K_P e = 0\) & Closed-loop error dynamics (perfect model) \\
\(V = \tfrac{1}{2}\dot{e}^T M \dot{e} + \tfrac{1}{2}e^T K_P e\) & Lyapunov function \\
\(\dot{V} = -\dot{e}^T K_D \dot{e} \leq 0\) & Lyapunov derivative (guarantees stability) \\
\end{longtable}

\subsection{Next Week: Advanced Control Techniques}

\textbf{Building on this week\textquotesingle s foundation, we will cover:}

\begin{itemize}
\tightlist
\item
  📐 Adaptive control --- automatically estimating \(M, C, G\) online
\item
  🛡️ Sliding mode control --- guaranteed robustness despite bounded uncertainties
\item
  🤖 Task-space (operational space) control
\item
  🔄 Impedance and force control for contact tasks
\end{itemize}

\subsection{Preparation for Next Lecture}

\begin{itemize}
\tightlist
\item
  Complete all MATLAB activities from today and save your scripts
\item
  Review Lyapunov stability theory (especially LaSalle\textquotesingle s invariance principle)
\item
  Read about sliding mode control fundamentals
\item
  Think about: \emph{"If I don\textquotesingle t know \(m_2\), how could the controller learn it during operation?"}
\end{itemize}

\textbf{✅ Skills Achieved Today:} You can now derive manipulator dynamics, design and implement a computed torque controller, prove its stability using Lyapunov theory, and analyse its robustness to model uncertainties --- the foundational skill set for all advanced robot control!

\textbf{Control for Robots}

Advanced Robotics Programme \textbar{} Academic Year 2025--26

Week 1: Manipulator Dynamics \& Control Design

This lecture material is for educational use only.
